\section{An atomic-time observation, literal latch and domain bindings}
\label{sec:r14-clock-binding}
An actual NIST internet-time observation is included as a fixed input to this
formalization. Its literal words provide a reproducible event anchor for the
physical models. They do not cause a running clock, calibrate an oscillator,
or supply the initial state of every physical system at that event.

The SI atomic reference is the unperturbed ground-state hyperfine transition
frequency of cesium-133:
\begin{equation}
 \Delta\nu_{\rm Cs}=9192631770\,\mathrm{Hz},\qquad
 1\,\mathrm{s}=9192631770/(\Delta\nu_{\rm Cs}) .
 \label{eq:r14-atomic-reference}
\end{equation}
The defining value is exact; a realization has measurement uncertainty.
\href{https://www.bipm.org/en/measurement-units/si-defining-constants}{[CL4: BIPM, SI defining constants]}
UTC(NIST) uses an ensemble of cesium-beam and hydrogen-maser clocks calibrated
against NIST's primary frequency standard. Its internet service distributes
that atomic reference; this capture does not install an atomic oscillator
in aTOMos.
\href{https://www.nist.gov/pml/time-and-frequency-division/how-utcnist-related-coordinated-universal-time-utc-international}{[CL5: NIST, UTC(NIST) realization]}

\subsection{The observed words and their exact rational coordinate}
The recorded exchange with \code{time.nist.gov}, endpoint
\code{132.163.96.4:123}, used a version-4 request and returned a version-3
response, mode 4, stratum 1, leap indicator 0 and reference ID \code{NIST}.
The complete request, reply and local acquisition stamps are preserved in
\path{source/clock_capture/nist_capture.json}. The server transmit field is
\begin{equation}
 w_{\rm tx}=17173213285462662115
 \quad\text{as an unsigned 64-bit word},\qquad e=0.
 \label{eq:r14-observed-ntp}
\end{equation}
The era $e$ is bound by the recorded contemporary-date context; the packet
does not carry it. The reply's SHA-256 is
\begin{center}\small
\path{26a429327733f35d8b67abfb9bbb33f54fae9fa590c4ba0e5a7947fe6946c3e4}.
\end{center}

For the 32-bit seconds and fraction fields $s,f\in\{0,\ldots,2^{32}-1\}$,
the mathematical decoder is
\begin{equation}
 w=2^{32}s+f,\qquad
 q_N=e\,2^{32}+s+\frac{f}{2^{32}},\qquad
 q_U=q_N-2208988800 .
 \label{eq:r14-ntp-rational}
\end{equation}
Here $q_N$ uses the 1900 NTP epoch and $q_U$ the corresponding 1970 civil
coordinate. Every operation is integer/rational; no floating conversion is
needed. The observed word gives the reduced rational
\begin{equation}
 q_U=\frac{7685678632232377315}{4294967296},
 \label{eq:r14-observed-rational}
\end{equation}
whose UTC display, truncated to microseconds, is
\code{2026-09-15T08:38:22.859456Z}. The display is not the authoritative value.
The NTP fixed-point format and era distinction follow
\href{https://www.rfc-editor.org/rfc/rfc5905.html#section-6}{[CL1: RFC 5905, section 6]}.

NTP's coordinate is affected by leap seconds; differences across leap events
are not automatically elapsed SI durations. A complete mapping binds the
era evidence, leap table, applicable interval and any ambiguity/smear policy.
\href{https://www.rfc-editor.org/rfc/rfc8877.html#section-4.2.1}{[CL2: RFC 8877, NTP timestamps]}
NIST documents its positive-leap convention as repetition of the final
second's binary timestamp.
\href{https://www.nist.gov/pml/time-and-frequency-division/time-distribution/internet-time-service-its}{[CL3: NIST Internet Time Service]}
This convention is part of the source identity; replacing the provider does
not silently inherit it. The earlier Cloudflare capture is retained separately
as acquisition provenance; the NIST record is the selected latch.

\Needspace{11\baselineskip}
\subsection{Complete nominal mapping at the captured epoch}
The retained \path{source/clock_capture/leap-seconds.list} has last transition
$q_N=3692217600$ with TAI--UTC offset $37$ seconds and declared expiry
$q_N=4023129600$. Its SHA-256 is
\begin{center}\small
\path{db5a895f16853b03bfc865e8d68f9fc8710ef1740e3400c701cd46a5bbbc3433}.
\end{center}
For the captured non-leap instant in this bound interval, the nominal
conversion body is explicitly
\begin{align}
 Q_{\rm TAI}&=q_U+37,\nonumber\\
 t_{\rm GPST}&=q_U-315964800+(37-19),\nonumber\\
 Q_{\rm TT}&=q_U+37+\frac{4023}{125},\nonumber\\
 Q_0&=220924800+\frac{4023}{125},\qquad
 t_{\rm TCG}=\frac{Q_{\rm TT}-Q_0}{1-L_G}.
 \label{eq:r14-captured-time-conversion}
\end{align}
GPST is elapsed seconds from 1980-01-06 00:00:00 GPST. $Q_{\rm TAI}$ and
$Q_{\rm TT}$ use 1970-01-01 00:00:00 calendar-coordinate origins on their
respective scales. $t_{\rm TCG}$ is elapsed TCG seconds from
1977-01-01 00:00:32.184, when TT and TCG are defined equal; $Q_0$ is that
event's TT coordinate and $L_G=6.969290134\times10^{-10}$ is exact.
These are finite rational expressions for the nominal scales. The leap file
is a dated dependency, not authority for an unannounced future leap or an
ambiguous repeated second. UTC(NIST)--UTC and GNSS realization differences
remain physical uncertainty.
\href{https://data.iana.org/time-zones/data/leap-seconds.list}{[CL6: IANA, retained leap table]}
\href{https://iers-conventions.obspm.fr/content/chapter10/tn36_c10.pdf}{[GT3: IERS, time coordinates]}

\subsection{What the clock observation establishes}
The timestamp denotes the server's transmit event. Its observed local
monotonic round-trip interval is the exact recorded count
$127594000/10^9$ seconds under the local timer convention. Neither that count
nor stratum, precision or dispersion metadata certifies a UTC error bound.
The UDP reply was not cryptographically authenticated. The host clock was
not changed and no oscillator calibration was performed.

Let $q_*$ be the nominal converted source-event coordinate and write
\begin{equation}
 t_{\rm tx}=q_*+\epsilon_s,\qquad
 t_{\rm rx}=t_{\rm tx}+d_{\leftarrow},\qquad d_{\leftarrow}\geq0 .
 \label{eq:r14-clock-event-error}
\end{equation}
Source clock error $\epsilon_s$, inbound delay $d_{\leftarrow}$, local readout
error and conversion/realization error enter one joint uncertainty set.
An interval for any of them needs supporting assumptions or calibration.
One exchange does not determine asymmetric one-way delays or prove source
authenticity. Literal retention preserves the observation actually received.

For a complete four-stamp exchange on compatible, unfolded coordinates,
the conventional offset and round-trip expressions are
\begin{equation}
 \widehat\theta=\frac{(T_2-T_1)+(T_3-T_4)}2,\qquad
 \widehat d=(T_4-T_1)-(T_3-T_2).
 \label{eq:r14-ntp-offset}
\end{equation}
These expressions can remain exact rationals. Interpreting
$\widehat\theta$ as the physical offset requires a clock/path model;
unequal inbound and outbound delays contribute an offset error.
\href{https://www.rfc-editor.org/rfc/rfc5905.html#section-8}{[CL1: RFC 5905, section 8]}

\subsection{Immutable XOP snapshot and a one-time latch}
Define a finite typed tuple $\mathcal C$ containing both raw packet words,
source identity, era and its evidence, acquisition stamps, protocol fields,
the rational coordinate and the complete time-mapping dependency identities.
The packet words use \code{Bits(384)} for this 48-octet capture. Other fields
use existing integer, rational, Boolean and fixed structural types; the
binding fixes their widths and ordering. No new timestamp opcode is needed.

The completed capture may be a closed dependency value. A newly acquired
capture instead enters through a sample port. Every consumer at a logical
step reads the same immutable old-state/input snapshot. The physical mapping
is an imported complete pure template, with its declared domain, rather than
an instruction to consult the network again.
The materialized \path{source/clock_capture/anchor_state.json} records the
selected NIST observation and nominal rational mappings. The equations here
specify its XOP binding; no running XOP clock executor is claimed.

Let $a_n$ be an accepted-record pulse, $v_n$ the admission predicate for its
typed record, $\ell_n$ the old latched Boolean, and $H_n$ the stored tuple.
The explicit one-time latch is
\begin{align}
 e_n&=\operatorname{IF}(\ell_n,\mathrm{false},
       \operatorname{IF}(a_n,v_n,\mathrm{false})),\nonumber\\
 \ell_{n+1}&=\ell_n\vee e_n,\qquad
 H_{n+1}=\operatorname{IF}(e_n,\mathcal C_n,H_n).
 \label{eq:r14-clock-latch}
\end{align}
Initially $\ell_0=\mathrm{false}$ and $H_0$ is a fixed closed placeholder of
the declared type, which has no observed-clock meaning while unlatched.
For defined inputs the pulse equals $a_n\wedge v_n\wedge\neg\ell_n$.
The displayed lazy form demands no new record, admission check or pulse once
latched, and demands no record when an unlatched acceptance pulse is false.
Both next roots commit together. Once $\ell_n$ is true, $e_n$ is false at
every later step and induction proves $H_{n+k}=H_n$ for all $k\geq0$.
Thus future internet replies cannot alter this seed's latched anchor.
If an electrical rising edge supplies acceptance, a declared old Boolean
$b_n$ gives $a_n=c_n\wedge\neg b_n$ and
$b_{n+1}=\operatorname{IF}(\ell_n,b_n,c_n)$ in the same commit.
Thus the optional edge memory also demands no new level once latched.
A later anchor is a separately versioned capture, not an implicit rewrite.

The raw sample time and canonical-time value remain distinct: R11 tag 5
exposes the source rational and tag 6 applies the complete \code{TimeConvention}
mapping to GPST-SI. An unresolved era or leap branch produces an explicit
non-value, not a guessed instant. Closing the captured record as a seed
dependency also closes its mapping table and source interval. No live lookup
is hidden in a pure expression.

\subsection{One anchor, explicit time maps for each physical domain}
Let $t_0$ denote the nominal latched event on the selected canonical scale.
Each domain uses a declared map $t_d=\chi_d(t)$ and its own state, frame,
boundary conditions and observation history. For a differentiable increasing
map, the equation transforms by the chain rule:
\begin{equation}
 \frac{dx_d}{dt}=\frac{d\chi_d}{dt}
 F_d\bigl(\chi_d(t),x_d,u_d\bigr),\qquad x_d(t_0)=x_{d,0}.
 \label{eq:r14-domain-time-map}
\end{equation}
The supplied state $x_{d,0}$ must refer to that event. An old observation
cannot become a current state merely by changing its timestamp.

\begin{longtable}{P{.24\textwidth}P{.68\textwidth}}
\toprule
Domain & Explicit binding to the latched event\\\midrule
\endhead
Electromagnetism and photonics & $E(x,t),B(x,t),P_j(x,t)$ and their source/boundary histories use the mapped field time. A carrier phase $\varphi(t)=\varphi_0+\int_{t_0}^{t}\omega(s)\,ds$ also needs its phase/frequency reference; NTP does not measure optical phase.\\
Mechanics and control & $q(t_0),p(t_0)$, attitude, loads and actuator events retain their initial-event identity. ASA/NA and JK transitions act through their explicit physical ports.\\
Thermal and material history & $T(t_0)$, internal variables and accumulated work/heat require the actual state/history binding. A clock latch does not reset energy or entropy.\\
Gravity and orbit & The GPST/TT/TCG mapping, compatible frame/force parameters and station/satellite events remain those of Section~\ref{r14-gravity-time}. A coordinate-time conversion does not measure proper clock time.\\
Sensors and calibration & Acquisition endpoints, integration windows, latency and calibration age use mapped recorded events, with the shared clock error in the observation uncertainty.\\
Digital state and replay & The integer transition index $n$ is causal order. A time $t_n=t_0+\sum_{j<n}h_j$ requires explicit durations $h_j$; $n$ alone is not seconds.\\
\bottomrule
\end{longtable}

A common event does not require equal sample rates or spatial rasterization.
Field expressions may retain their continuous variables and exact implicit
denotation. A selected finite mode basis or finite step remains a separately
specified approximation with its own discrepancy. Clock anchoring supplies
the temporal relationship; it does not turn a finite model into every
physical domain or remove source, state or model uncertainty.
